\subsection{Kỹ năng và kiến thức thu nhận được}

Trong kỳ thực tập tại \textbf{Công ty TNHH Amazon Web Services Việt Nam}, từ
\textbf{15/06/2026} đến \textbf{14/08/2026}, thuộc chương trình Workforce
Bootcamp -- First Cloud Journey, tôi đã chuyển từ việc học điện toán đám mây
trên lý thuyết sang thiết kế, xây dựng, và vận hành một hệ thống hoàn chỉnh
trên AWS.

Sáu tuần đầu được dành để đi qua các nhóm dịch vụ cốt lõi -- quản lý danh tính
và truy cập, mạng, compute, lưu trữ, cơ sở dữ liệu, serverless, và giám sát --
mỗi nhóm đều kèm theo một bài thực hành. Những tuần cuối được dành để áp dụng
toàn bộ những điều đó vào \textbf{Caerus}, một nền tảng đặt ghế xem phim do một
nhóm hai người xây dựng và triển khai trên Amazon EC2, Amazon RDS, Amazon S3,
API Gateway và Amazon CloudWatch.

Những kỹ năng cụ thể tôi thu nhận được từ chương trình gồm:

\begin{itemize}
    \setlength\itemsep{0.2em}
    \item \textbf{Kiến trúc đám mây và lựa chọn dịch vụ.} Tôi đã có thể thiết
    kế một kiến trúc nhiều dịch vụ và biện luận cho từng lựa chọn dịch vụ bằng
    cách đặt nó cạnh một phương án thay thế cụ thể, thay vì chọn theo mặc định.

    \item \textbf{Thiết kế cơ sở dữ liệu quan hệ, lập trình giao dịch và kiểm
    soát concurrency.} Đặt ghế trong điều kiện concurrency hoá ra là một bài
    toán giao dịch chứ không phải bài toán giao diện hay hạ tầng: khóa tường
    minh những dòng đang bị tranh chấp, lấy khóa theo một thứ tự nhất quán để
    tránh deadlock, rồi chứng minh rằng cam kết đó thực sự đúng.

    \item \textbf{Thiết kế API và làm việc dựa trên một bản hợp đồng.} Việc
    thống nhất bản đặc tả API và schema cơ sở dữ liệu trước khi viết bất kỳ dòng
    code nào chính là điều cho phép hai người cùng xây dựng song song thay vì
    người này phải chờ người kia.

    \item \textbf{Triển khai và cấu hình bảo mật mạng}, bao gồm security group,
    việc đặt tài nguyên vào đúng subnet, và định tuyến kết nối giữa chúng.

    \item \textbf{Khả năng quan sát và quản lý chi phí.} Dashboard và alarm cho
    vế thứ nhất; billing alarm, tag chủ sở hữu, và terminate tài nguyên thực
    hành ngay trong ngày cho vế thứ hai.

    \item \textbf{Viết tài liệu kỹ thuật.} Ghi lại các quyết định kiến trúc và
    quy trình vận hành sao cho một người không có mặt lúc đó vẫn theo được.
\end{itemize}

Về tinh thần làm việc, tôi bám sát lịch trình đã thống nhất từ đầu dự án, duy
trì những thói quen hằng ngày mà nhóm đã cam kết -- terminate tài nguyên thực
hành ngay trong ngày và gắn tag chủ sở hữu cho mọi tài nguyên -- và nêu vấn đề
với đồng đội từ sớm thay vì tự mình xoay xở vòng qua chúng.

\subsection{Những điều đã làm tốt}

\begin{itemize}
    \setlength\itemsep{0.2em}
    \item \textbf{Độ rộng của việc học.} Tôi bắt đầu chương trình mà không có
    kinh nghiệm thực tế nào với AWS và kết thúc nó với khả năng thiết kế một
    kiến trúc nhiều dịch vụ và bảo vệ được từng lựa chọn dịch vụ.

    \item \textbf{Làm việc dựa trên một bản hợp đồng.} Việc đóng băng bản đặc tả
    API và schema cơ sở dữ liệu ngay từ đầu là quyết định duy nhất giữ cho dự án
    đúng tiến độ. Khâu tích hợp trở thành chuyện sửa vài điểm lệch nhỏ thay vì
    hòa giải hai thiết kế không tương thích.

    \item \textbf{Giải đúng bài toán mà dự án thực sự nói về.} Tôi kiểm chứng
    cam kết về concurrency bằng một cuộc chạy đua hai client có chủ đích thay vì
    mặc định rằng nó đúng.

    \item \textbf{Kỷ luật về chi phí.} Việc đặt một billing alarm trước khi khởi
    tạo bất cứ thứ gì, gắn tag chủ sở hữu cho mọi tài nguyên, và terminate tài
    nguyên thực hành ngay trong ngày đã khiến chi phí không bao giờ trở thành
    một vấn đề phải xử lý về sau.
\end{itemize}

\subsection{Những điều cần cải thiện}

\begin{itemize}
    \setlength\itemsep{0.2em}
    \item \textbf{Chủ động vượt ra ngoài kế hoạch.} Tôi thực thi lịch trình đã
    thống nhất một cách đáng tin cậy, nhưng lại có xu hướng làm hết danh sách
    công việc thay vì nhìn xa hơn và đề xuất cải tiến trước khi chúng trở thành
    bắt buộc. Nhiều vấn đề mà nhóm gặp trong lúc triển khai vốn có thể lường
    trước, và lẽ ra tôi nên nêu chúng ra ngay ở khâu thiết kế thay vì phát hiện
    ra khi đang chịu áp lực thời gian.

    \item \textbf{Giao tiếp và thuyết trình.} Tôi thoải mái khi giải thích các
    quyết định kỹ thuật bằng văn bản, nhưng kém tự tin hơn khi trình bày chúng
    bằng lời trước một người chưa đọc tài liệu. Đây là kỹ năng tôi muốn phát
    triển nhất trong thời gian tới.

    \item \textbf{Chiều sâu vượt ra ngoài con đường chạy được.} Việc triển khai
    vẫn làm tay chứ chưa tự động hóa, hạ tầng được tạo qua console chứ chưa định
    nghĩa dưới dạng code, và kiểm thử tự động mới giới hạn ở trường hợp
    concurrency chứ chưa phủ rộng ra toàn bộ API.

    \item \textbf{Kiến trúc ở mức production.} Ứng dụng chạy trên một instance
    duy nhất trong một public subnet. Tôi hiểu một bản triển khai vững vàng hơn
    sẽ cần những gì -- private subnet với truy cập theo phiên, một load
    balancer, một auto scaling group, và một cơ sở dữ liệu multi-AZ -- nhưng tôi
    chưa từng tự tay dựng lên, nên hiểu biết của tôi về những mẫu đó vẫn còn ở
    mức khái niệm chứ chưa phải thực hành.

    \item \textbf{Ước lượng thời gian.} Các ước lượng ban đầu của tôi cho khâu
    triển khai và cấu hình cross-origin đều lạc quan. Những ngày dự phòng cài
    sẵn trong kế hoạch đã hấp thụ được phần vượt tiến độ, nhưng bản thân khả
    năng ước lượng thì vẫn cần cải thiện.
\end{itemize}

\subsection{Tự đánh giá}

Để nhìn lại kỳ thực tập một cách khách quan, tôi tự đánh giá theo các tiêu chí
sau:

\begin{center}
\small
\renewcommand{\arraystretch}{1.25}
\begin{tabular}{@{}c>{\raggedright\arraybackslash}p{7.2cm}ccc@{}}
\toprule
\textbf{STT} & \textbf{Tiêu chí} & \textbf{Tốt} & \textbf{Khá} & \textbf{TB} \\
\midrule
1  & Kiến thức \& kỹ năng chuyên môn & $\checkmark$ & $\square$ & $\square$ \\
2  & Khả năng học hỏi                & $\checkmark$ & $\square$ & $\square$ \\
3  & Tính chủ động                   & $\square$ & $\checkmark$ & $\square$ \\
4  & Tinh thần trách nhiệm           & $\checkmark$ & $\square$ & $\square$ \\
5  & Tính kỷ luật                    & $\checkmark$ & $\square$ & $\square$ \\
6  & Tinh thần cầu tiến              & $\checkmark$ & $\square$ & $\square$ \\
7  & Giao tiếp                       & $\square$ & $\checkmark$ & $\square$ \\
8  & Làm việc nhóm                   & $\checkmark$ & $\square$ & $\square$ \\
9  & Tác phong chuyên nghiệp         & $\checkmark$ & $\square$ & $\square$ \\
10 & Kỹ năng giải quyết vấn đề       & $\checkmark$ & $\square$ & $\square$ \\
11 & Đóng góp cho dự án/nhóm         & $\checkmark$ & $\square$ & $\square$ \\
12 & Đánh giá tổng thể               & $\checkmark$ & $\square$ & $\square$ \\
\bottomrule
\end{tabular}
\end{center}

\subsection{Cảm nhận cá nhân}

Chính cách tổ chức của chương trình là điều làm cho bước nhảy từ lý thuyết sang
thực hành trở nên khả thi. Sáu tuần đi qua các nhóm dịch vụ cốt lõi, mỗi nhóm
kèm một bài thực hành, cho tôi đủ nền tảng để những tuần cuối được dành cho một
hệ thống thật thay vì dành để học cách dùng console. Việc được yêu cầu xây dựng
một sản phẩm từ đầu đến cuối, và phải biện luận cho kiến trúc chứ không chỉ lắp
ghép nó, hữu ích hơn nhiều so với bất kỳ bài hướng dẫn riêng lẻ nào về một dịch
vụ.

Điều tôi không ngờ mình lại coi trọng đến vậy là kỷ luật vận hành mà chương
trình yêu cầu ngay từ ngày đầu: đặt billing alarm trước khi khởi tạo bất cứ thứ
gì, gắn tag chủ sở hữu cho mọi tài nguyên, và terminate tài nguyên thực hành
ngay trong ngày. Những thói quen đó gần như không tốn công duy trì nhưng đã loại
bỏ hẳn một nhóm vấn đề khỏi dự án. Đây là phần tôi tin mình sẽ mang theo vào mọi
dự án về sau.

Việc làm việc theo cặp dạy cho tôi phần còn lại. Thống nhất một hợp đồng giao
diện trước khi viết code không phải là thủ tục rườm rà; đó chính là điều cho
phép hai người xây dựng song song. Tôi xin cảm ơn cán bộ hướng dẫn tại doanh
nghiệp và cán bộ hướng dẫn của khoa đã đồng hành cùng chúng tôi trong suốt kỳ
thực tập.

% Phần chữ ký được đặt riêng một trang, canh ở đầu trang
% (dùng \vspace* thay vì \vspace, vì \vspace bị bỏ qua ngay sau khi ngắt trang).
\newpage
\vspace*{1cm}

% Giữ ba ô chữ ký trên cùng một hàng để không bị ngắt sang trang khác nhau.
\noindent
\begin{minipage}[t]{0.31\textwidth}
    \centering
    \textbf{Sinh viên}\\[0.2cm]
    % (Kí và ghi rõ họ tên)

    \vspace{2.2cm}

    \rule{0.9\linewidth}{0.4pt}
\end{minipage}
\hfill
\begin{minipage}[t]{0.31\textwidth}
    \centering
    \textbf{Sinh viên}\\[0.2cm]
    % (Kí và ghi rõ họ tên)

    \vspace{2.2cm}

    \rule{0.9\linewidth}{0.4pt}
\end{minipage}
\hfill
\begin{minipage}[t]{0.31\textwidth}
    \centering
    \textbf{Người đại diện}\\\textbf{doanh nghiệp}\\[0.2cm]
    % (Kí và ghi rõ họ tên)

    \vspace{1.85cm}

    \rule{0.9\linewidth}{0.4pt}
\end{minipage}
